\begin{PSbox}[unbreakable]{obj}{Learning Objectives}

After completing this module, students will be able to:

\begin{enumerate}
\def\labelenumi{\arabic{enumi}.}
\tightlist
\item
  Compute conditional expectations and variances
\item
  Apply the Law of Total Expectation
\item
  Apply the Law of Total Variance
\item
  Use these tools for engineering analysis across multiple operating conditions
\end{enumerate}

\end{PSbox}

\PSsection{8.1}{Introduction}

Many engineering systems operate under \textbf{different conditions} (channel states, modes, environments). We often need the \textbf{overall average performance} computed from conditional performance in each state.

\begin{itemize}
\tightlist
\item
  What is the average BER when the channel randomly switches between good and bad?
\item
  What is the overall measurement variance when the sensor operates in different environments?
\end{itemize}

\PSsection{8.2}{Conditional Expectation $E[X \mid Y = y]$}

\begin{PSbox}[unbreakable]{def}{Definition}

\textbf{Continuous:}\PSdisplay{E[X|Y=y] = \int_{-\infty}^{\infty} x \cdot f_{X|Y}(x|y) \, dx}

\textbf{Discrete:}\PSdisplay{E[X|Y=y] = \sum_x x \cdot p_{X|Y}(x|y)}

\end{PSbox}

\PSsubsection{Interpretation}

$E[X \mid Y = y]$ is the \textbf{average value of $X$} when we know $Y = y$. It is a function of $y$.

\PSsubsection{$E[X \mid Y]$ as a Random Variable}

When $Y$ is itself random, $E[X \mid Y]$ is a \textbf{random variable} — it takes the value $E[X \mid Y = y]$ when $Y = y$.

\begin{PSbox}[unbreakable]{ex}{Engineering Example}

Average received power given channel state:

\begin{itemize}
\tightlist
\item
  Good channel $(Y = 1)$: $E[P \mid Y = 1] = 10$ mW
\item
  Poor channel $(Y = 0)$: $E[P \mid Y = 0] = 1$ mW
\item
  $P(Y = 1) = 0.7,\, P(Y = 0) = 0.3$
\end{itemize}

\end{PSbox}

\PSsection{8.3}{Conditional Variance $\operatorname{Var}(X \mid Y = y)$}

\begin{PSbox}[unbreakable]{def}{Definition}

\PSdisplay{\text{Var}(X|Y=y) = E[(X - E[X|Y=y])^2 | Y=y] = E[X^2|Y=y] - (E[X|Y=y])^2}

\end{PSbox}

\PSsubsection{Interpretation}

The conditional variance measures the \textbf{remaining uncertainty in $X$} after $Y$ is known.

\begin{PSbox}[unbreakable]{ex}{Engineering Example}

Measurement precision depends on temperature:

\begin{itemize}
\tightlist
\item
  At $T = 20^{\circ}C$: $\operatorname{Var}(X \mid T = 20) = 0.01$ (high precision)
\item
  At $T = 80^{\circ}C$: $\operatorname{Var}(X \mid T = 80) = 0.09$ (low precision due to thermal effects)
\end{itemize}

\end{PSbox}

\PSsection{8.4}{Law of Total Expectation (Tower Property)}

\PSsubsection{Statement}

\PSdisplay{E[X] = E[E[X|Y]]}

\textbf{Discrete version:}\PSdisplay{E[X] = \sum_y E[X|Y=y] \cdot P(Y=y)}

\textbf{Continuous version:}\PSdisplay{E[X] = \int_{-\infty}^{\infty} E[X|Y=y] \cdot f_Y(y) \, dy}

\PSsubsection{Proof (Discrete)}

$E[E[X \mid Y]] = \sum_{y} E[X \mid Y = y] \cdot P(Y = y)$\PSbr{}$= \sum_{y} \left[\sum_{x} x \cdot P(X = x \mid Y = y)\right] \cdot P(Y = y)$\PSbr{}$= \sum_{x} \Sigma_{y} x \cdot P(X = x,\, Y = y)$\PSbr{}$= \sum_{x} x \cdot P(X = x) = E[X]$ \PSyes{}

\PSsubsection{Interpretation}

The \textbf{overall average} equals the \textbf{weighted average of conditional averages}, weighted by the probability of each condition.

\begin{PSbox}[unbreakable]{ex}{Engineering Example: Average BER Across Channel States}

A wireless system switches between three modes:

\begin{itemize}
\tightlist
\item
  Mode 1 (60\% of time): $\mathrm{BER} = 10^{-6}$
\item
  Mode 2 (30\% of time): $\mathrm{BER} = 10^{-4}$
\item
  Mode 3 (10\% of time): $\mathrm{BER} = 10^{-2}$
\end{itemize}

Overall $\mathrm{BER} = E[\mathrm{BER}] = 0.6(10^{-6}) + 0.3(10^{-4}) + 0.1(10^{-2})$\PSbr{}$= 0.6 \times 10^{-6} + 30 \times 10^{-6} + 10000 \times 10^{-6} \approx 1.003 \times 10^{-3}$

\textbf{Engineering insight:} The overall BER is dominated by the worst mode, even though it occurs only 10\% of the time.

\end{PSbox}

\begin{PSbox}[unbreakable]{ex}{Engineering Example: Average Lifetime with Defect Types}

Components are defect-free (90\%) or have a hidden defect (10\%):

\begin{itemize}
\tightlist
\item
  Defect-free: $E[T \mid \text{good}] = 10000$ hours
\item
  With defect: $E[T \mid \text{defect}] = 500$ hours
\end{itemize}

$E[T] = 0.9(10000) + 0.1(500) = 9050$ hours

\end{PSbox}

\PSsection{8.5}{Law of Total Variance}

\PSsubsection{Statement}

\PSdisplay{\text{Var}(X) = E[\text{Var}(X|Y)] + \text{Var}(E[X|Y])}

{\def\LTcaptype{none} % do not increment counter
\begin{longtable}[c]{@{}
  >{\centering\arraybackslash}p{(0.965\linewidth - 4\tabcolsep) * \real{0.2222}}
  >{\centering\arraybackslash}p{(0.965\linewidth - 4\tabcolsep) * \real{0.2222}}
  >{\centering\arraybackslash}p{(0.965\linewidth - 4\tabcolsep) * \real{0.5556}}@{}}
\toprule\noalign{}
\rowcolor{PSnavy}\PSth{Term} & \PSth{Name} & \PSth{Interpretation} \\
\midrule\noalign{}
\endhead
\bottomrule\noalign{}
\endlastfoot
$\operatorname{Var}(X)$ & Total variance & Overall uncertainty in $X$ \\
$E[\operatorname{Var}(X \mid Y)]$ & Expected conditional variance & Average ``within-group'' variability \\
$\operatorname{Var}(E[X \mid Y])$ & Variance of conditional means & ``Between-group'' variability \\
\end{longtable}
}

\PSsubsection{Proof}

$\operatorname{Var}(X) = E[X^{2}] - (E[X])^{2}$

$E[X^{2}] = E[E[X^{2} \mid Y]]$ (tower property)\PSbr{}$= E[\operatorname{Var}(X \mid Y) + (E[X \mid Y])^{2}]$\PSbr{}$= E[\operatorname{Var}(X \mid Y)] + E[(E[X \mid Y])^{2}]$

$(E[X])^{2} = (E[E[X \mid Y]])^{2}$

$\operatorname{Var}(X) = E[\operatorname{Var}(X \mid Y)] + E[(E[X \mid Y])^{2}] - (E[E[X \mid Y]])^{2}$\PSbr{}$= E[\operatorname{Var}(X \mid Y)] + \operatorname{Var}(E[X \mid Y])$ \PSyes{}

\PSsubsection{Interpretation: Decomposition of Uncertainty}

Total uncertainty = Average uncertainty within each condition + Variability of the mean across conditions

\begin{PSbox}[unbreakable]{ex}{Engineering Example: Sensor in Two Environments}

A sensor operates indoors ($Y = 1$, prob 0.6) and outdoors ($Y = 2$, prob 0.4):

\begin{itemize}
\tightlist
\item
  Indoor: $E[X \mid Y = 1] = 25^{\circ}C,\, \operatorname{Var}(X \mid Y = 1) = 1$
\item
  Outdoor: $E[X \mid Y = 2] = 15^{\circ}C,\, \operatorname{Var}(X \mid Y = 2) = 16$
\end{itemize}

**$E[\operatorname{Var}(X \mid Y)] \ast \ast = 0.6(1) + 0.4(16) = 7.0$ (average within-condition variance)

**$\operatorname{Var}(E[X \mid Y]) \ast \ast = E[(E[X \mid Y])^{2}] - (E[E[X \mid Y]])^{2}$\PSbr{}$= 0.6(25^{2}) + 0.4(15^{2}) - [0.6(25) + 0.4(15)]^{2} = 375 + 90 - (21)^{2} = 465 - 441 = 24$

**$\operatorname{Var}(X) \ast \ast = 7.0 + 24.0 = 31.0$

Most variance comes from the \textbf{between-condition} difference (24 out of 31).

\end{PSbox}

\PSsection{8.6}{MATLAB Examples}

\begin{PSbox}[unbreakable]{ex}{Example 1: Law of Total Expectation}

\tcbset{PScodemode/.style={unbreakable}}

\begin{Shaded}
\begin{Highlighting}[]
\CommentTok{\%\% Average BER across channel states}
\CommentTok{\% State probabilities and conditional BERs}
\VariableTok{p\_state} \OperatorTok{=}\NormalTok{ [}\FloatTok{0.6}\OperatorTok{,} \FloatTok{0.3}\OperatorTok{,} \FloatTok{0.1}\NormalTok{]}\OperatorTok{;}
\VariableTok{BER\_cond} \OperatorTok{=}\NormalTok{ [}\FloatTok{1e{-}6}\OperatorTok{,} \FloatTok{1e{-}4}\OperatorTok{,} \FloatTok{1e{-}2}\NormalTok{]}\OperatorTok{;}

\CommentTok{\% Total expectation}
\VariableTok{BER\_overall} \OperatorTok{=} \VariableTok{sum}\NormalTok{(}\VariableTok{p\_state} \OperatorTok{.*} \VariableTok{BER\_cond}\NormalTok{)}\OperatorTok{;}
\VariableTok{fprintf}\NormalTok{(}\SpecialStringTok{\textquotesingle{}Overall BER = \%.4e\textbackslash{}n\textquotesingle{}}\OperatorTok{,} \VariableTok{BER\_overall}\NormalTok{)}\OperatorTok{;}

\CommentTok{\% Simulation verification}
\VariableTok{N} \OperatorTok{=} \FloatTok{1000000}\OperatorTok{;}
\VariableTok{states} \OperatorTok{=} \VariableTok{randsample}\NormalTok{(}\FloatTok{1}\OperatorTok{:}\FloatTok{3}\OperatorTok{,} \VariableTok{N}\OperatorTok{,} \VariableTok{true}\OperatorTok{,} \VariableTok{p\_state}\NormalTok{)}\OperatorTok{;}
\VariableTok{errors} \OperatorTok{=} \VariableTok{rand}\NormalTok{(}\FloatTok{1}\OperatorTok{,}\VariableTok{N}\NormalTok{) }\OperatorTok{\textless{}} \VariableTok{BER\_cond}\NormalTok{(}\VariableTok{states}\NormalTok{)}\OperatorTok{;}
\VariableTok{BER\_sim} \OperatorTok{=} \VariableTok{mean}\NormalTok{(}\VariableTok{errors}\NormalTok{)}\OperatorTok{;}
\VariableTok{fprintf}\NormalTok{(}\SpecialStringTok{\textquotesingle{}Simulated BER = \%.4e\textbackslash{}n\textquotesingle{}}\OperatorTok{,} \VariableTok{BER\_sim}\NormalTok{)}\OperatorTok{;}
\end{Highlighting}
\end{Shaded}

\end{PSbox}

\begin{PSbox}[unbreakable]{ex}{Example 2: Law of Total Variance}

\tcbset{PScodemode/.style={unbreakable}}

\begin{Shaded}
\begin{Highlighting}[]
\CommentTok{\%\% Total Variance Decomposition: Sensor Example}
\VariableTok{N} \OperatorTok{=} \FloatTok{100000}\OperatorTok{;}
\VariableTok{p\_indoor} \OperatorTok{=} \FloatTok{0.6}\OperatorTok{;}

\CommentTok{\% Generate environment (1=indoor, 2=outdoor)}
\VariableTok{env} \OperatorTok{=}\NormalTok{ (}\VariableTok{rand}\NormalTok{(}\FloatTok{1}\OperatorTok{,}\VariableTok{N}\NormalTok{) }\OperatorTok{\textgreater{}} \VariableTok{p\_indoor}\NormalTok{) }\OperatorTok{+} \FloatTok{1}\OperatorTok{;}

\CommentTok{\% Generate measurements conditional on environment}
\VariableTok{X} \OperatorTok{=} \VariableTok{zeros}\NormalTok{(}\FloatTok{1}\OperatorTok{,} \VariableTok{N}\NormalTok{)}\OperatorTok{;}
\VariableTok{X}\NormalTok{(}\VariableTok{env}\OperatorTok{==}\FloatTok{1}\NormalTok{) }\OperatorTok{=} \FloatTok{25} \OperatorTok{+} \FloatTok{1}\OperatorTok{*}\VariableTok{randn}\NormalTok{(}\FloatTok{1}\OperatorTok{,} \VariableTok{sum}\NormalTok{(}\VariableTok{env}\OperatorTok{==}\FloatTok{1}\NormalTok{))}\OperatorTok{;}    \CommentTok{\% Indoor: N(25, 1)}
\VariableTok{X}\NormalTok{(}\VariableTok{env}\OperatorTok{==}\FloatTok{2}\NormalTok{) }\OperatorTok{=} \FloatTok{15} \OperatorTok{+} \FloatTok{4}\OperatorTok{*}\VariableTok{randn}\NormalTok{(}\FloatTok{1}\OperatorTok{,} \VariableTok{sum}\NormalTok{(}\VariableTok{env}\OperatorTok{==}\FloatTok{2}\NormalTok{))}\OperatorTok{;}    \CommentTok{\% Outdoor: N(15, 16)}

\CommentTok{\% Total variance}
\VariableTok{total\_var} \OperatorTok{=} \VariableTok{var}\NormalTok{(}\VariableTok{X}\NormalTok{)}\OperatorTok{;}

\CommentTok{\% E[Var(X|Y)]}
\VariableTok{within\_var} \OperatorTok{=} \VariableTok{p\_indoor}\OperatorTok{*}\VariableTok{var}\NormalTok{(}\VariableTok{X}\NormalTok{(}\VariableTok{env}\OperatorTok{==}\FloatTok{1}\NormalTok{)) }\OperatorTok{+}\NormalTok{ (}\FloatTok{1}\OperatorTok{{-}}\VariableTok{p\_indoor}\NormalTok{)}\OperatorTok{*}\VariableTok{var}\NormalTok{(}\VariableTok{X}\NormalTok{(}\VariableTok{env}\OperatorTok{==}\FloatTok{2}\NormalTok{))}\OperatorTok{;}

\CommentTok{\% Var(E[X|Y])}
\VariableTok{cond\_means} \OperatorTok{=}\NormalTok{ [}\VariableTok{mean}\NormalTok{(}\VariableTok{X}\NormalTok{(}\VariableTok{env}\OperatorTok{==}\FloatTok{1}\NormalTok{))}\OperatorTok{,} \VariableTok{mean}\NormalTok{(}\VariableTok{X}\NormalTok{(}\VariableTok{env}\OperatorTok{==}\FloatTok{2}\NormalTok{))]}\OperatorTok{;}
\VariableTok{between\_var} \OperatorTok{=} \VariableTok{p\_indoor}\OperatorTok{*}\VariableTok{cond\_means}\NormalTok{(}\FloatTok{1}\NormalTok{)}\OperatorTok{\^{}}\FloatTok{2} \OperatorTok{+}\NormalTok{ (}\FloatTok{1}\OperatorTok{{-}}\VariableTok{p\_indoor}\NormalTok{)}\OperatorTok{*}\VariableTok{cond\_means}\NormalTok{(}\FloatTok{2}\NormalTok{)}\OperatorTok{\^{}}\FloatTok{2} \OperatorTok{...}
              \OperatorTok{{-}}\NormalTok{ (}\VariableTok{p\_indoor}\OperatorTok{*}\VariableTok{cond\_means}\NormalTok{(}\FloatTok{1}\NormalTok{) }\OperatorTok{+}\NormalTok{ (}\FloatTok{1}\OperatorTok{{-}}\VariableTok{p\_indoor}\NormalTok{)}\OperatorTok{*}\VariableTok{cond\_means}\NormalTok{(}\FloatTok{2}\NormalTok{))}\OperatorTok{\^{}}\FloatTok{2}\OperatorTok{;}

\VariableTok{fprintf}\NormalTok{(}\SpecialStringTok{\textquotesingle{}Total Var(X) = \%.2f (theory: 31.0)\textbackslash{}n\textquotesingle{}}\OperatorTok{,} \VariableTok{total\_var}\NormalTok{)}\OperatorTok{;}
\VariableTok{fprintf}\NormalTok{(}\SpecialStringTok{\textquotesingle{}E[Var(X|Y)]  = \%.2f (theory: 7.0)\textbackslash{}n\textquotesingle{}}\OperatorTok{,} \VariableTok{within\_var}\NormalTok{)}\OperatorTok{;}
\VariableTok{fprintf}\NormalTok{(}\SpecialStringTok{\textquotesingle{}Var(E[X|Y])  = \%.2f (theory: 24.0)\textbackslash{}n\textquotesingle{}}\OperatorTok{,} \VariableTok{between\_var}\NormalTok{)}\OperatorTok{;}
\VariableTok{fprintf}\NormalTok{(}\SpecialStringTok{\textquotesingle{}Sum           = \%.2f\textbackslash{}n\textquotesingle{}}\OperatorTok{,} \VariableTok{within\_var} \OperatorTok{+} \VariableTok{between\_var}\NormalTok{)}\OperatorTok{;}
\end{Highlighting}
\end{Shaded}

\end{PSbox}

\begin{PSbox}[unbreakable]{ex}{Example 3: Conditional Expectation as Best Predictor}

\tcbset{PScodemode/.style={unbreakable}}

\begin{Shaded}
\begin{Highlighting}[]
\CommentTok{\%\% E[X|Y] as a function of Y}
\CommentTok{\% Joint: X = Y + N, Y \textasciitilde{} Uniform(0,10), N \textasciitilde{} N(0, 1)}
\VariableTok{N\_samp} \OperatorTok{=} \FloatTok{50000}\OperatorTok{;}
\VariableTok{Y} \OperatorTok{=} \FloatTok{10}\OperatorTok{*}\VariableTok{rand}\NormalTok{(}\FloatTok{1}\OperatorTok{,} \VariableTok{N\_samp}\NormalTok{)}\OperatorTok{;}
\VariableTok{X} \OperatorTok{=} \VariableTok{Y} \OperatorTok{+} \VariableTok{randn}\NormalTok{(}\FloatTok{1}\OperatorTok{,} \VariableTok{N\_samp}\NormalTok{)}\OperatorTok{;}

\CommentTok{\% Estimate E[X|Y=y] for various y}
\VariableTok{y\_bins} \OperatorTok{=} \FloatTok{0}\OperatorTok{:}\FloatTok{0.5}\OperatorTok{:}\FloatTok{10}\OperatorTok{;}
\VariableTok{E\_X\_given\_Y} \OperatorTok{=} \VariableTok{zeros}\NormalTok{(}\VariableTok{size}\NormalTok{(}\VariableTok{y\_bins}\NormalTok{)}\OperatorTok{{-}}\NormalTok{[}\FloatTok{0} \FloatTok{1}\NormalTok{])}\OperatorTok{;}
\KeywordTok{for} \VariableTok{i} \OperatorTok{=} \FloatTok{1}\OperatorTok{:}\VariableTok{length}\NormalTok{(}\VariableTok{y\_bins}\NormalTok{)}\OperatorTok{{-}}\FloatTok{1}
    \VariableTok{mask} \OperatorTok{=} \VariableTok{Y} \OperatorTok{\textgreater{}=} \VariableTok{y\_bins}\NormalTok{(}\VariableTok{i}\NormalTok{) }\OperatorTok{\&} \VariableTok{Y} \OperatorTok{\textless{}} \VariableTok{y\_bins}\NormalTok{(}\VariableTok{i}\OperatorTok{+}\FloatTok{1}\NormalTok{)}\OperatorTok{;}
    \VariableTok{E\_X\_given\_Y}\NormalTok{(}\VariableTok{i}\NormalTok{) }\OperatorTok{=} \VariableTok{mean}\NormalTok{(}\VariableTok{X}\NormalTok{(}\VariableTok{mask}\NormalTok{))}\OperatorTok{;}
\KeywordTok{end}
\VariableTok{y\_centers} \OperatorTok{=}\NormalTok{ (}\VariableTok{y\_bins}\NormalTok{(}\FloatTok{1}\OperatorTok{:}\KeywordTok{end}\OperatorTok{{-}}\FloatTok{1}\NormalTok{) }\OperatorTok{+} \VariableTok{y\_bins}\NormalTok{(}\FloatTok{2}\OperatorTok{:}\KeywordTok{end}\NormalTok{))}\OperatorTok{/}\FloatTok{2}\OperatorTok{;}

\VariableTok{figure}\OperatorTok{;}
\VariableTok{scatter}\NormalTok{(}\VariableTok{Y}\NormalTok{(}\FloatTok{1}\OperatorTok{:}\FloatTok{5000}\NormalTok{)}\OperatorTok{,} \VariableTok{X}\NormalTok{(}\FloatTok{1}\OperatorTok{:}\FloatTok{5000}\NormalTok{)}\OperatorTok{,} \FloatTok{1}\OperatorTok{,} \SpecialStringTok{\textquotesingle{}b\textquotesingle{}}\OperatorTok{,} \SpecialStringTok{\textquotesingle{}.\textquotesingle{}}\NormalTok{)}\OperatorTok{;}
\VariableTok{hold} \VariableTok{on}\OperatorTok{;}
\VariableTok{plot}\NormalTok{(}\VariableTok{y\_centers}\OperatorTok{,} \VariableTok{E\_X\_given\_Y}\OperatorTok{,} \SpecialStringTok{\textquotesingle{}r{-}\textquotesingle{}}\OperatorTok{,} \SpecialStringTok{\textquotesingle{}LineWidth\textquotesingle{}}\OperatorTok{,} \FloatTok{3}\NormalTok{)}\OperatorTok{;}
\VariableTok{plot}\NormalTok{(}\VariableTok{y\_centers}\OperatorTok{,} \VariableTok{y\_centers}\OperatorTok{,} \SpecialStringTok{\textquotesingle{}g{-}{-}\textquotesingle{}}\OperatorTok{,} \SpecialStringTok{\textquotesingle{}LineWidth\textquotesingle{}}\OperatorTok{,} \FloatTok{2}\NormalTok{)}\OperatorTok{;}
\VariableTok{xlabel}\NormalTok{(}\SpecialStringTok{\textquotesingle{}Y\textquotesingle{}}\NormalTok{)}\OperatorTok{;} \VariableTok{ylabel}\NormalTok{(}\SpecialStringTok{\textquotesingle{}X\textquotesingle{}}\NormalTok{)}\OperatorTok{;}
\VariableTok{legend}\NormalTok{(}\SpecialStringTok{\textquotesingle{}Data (X,Y)\textquotesingle{}}\OperatorTok{,} \SpecialStringTok{\textquotesingle{}E[X|Y] (estimated)\textquotesingle{}}\OperatorTok{,} \SpecialStringTok{\textquotesingle{}E[X|Y] = Y (theory)\textquotesingle{}}\NormalTok{)}\OperatorTok{;}
\VariableTok{title}\NormalTok{(}\SpecialStringTok{\textquotesingle{}Conditional Expectation\textquotesingle{}}\NormalTok{)}\OperatorTok{;}
\end{Highlighting}
\end{Shaded}

\end{PSbox}

\PSsection{8.7}{Practice Problems}

\begin{PSbox}[unbreakable]{prob}{Problem 1}

A digital link operates over two channel types:

\begin{itemize}
\tightlist
\item
  Type $A$ (prob 0.7): $E[\text{delay}\ \mid A] = 5\,\mathrm{ms},\, \operatorname{Var}(\text{delay}\ \mid A) = 2$
\item
  Type $B$ (prob 0.3): $E[\text{delay}\ \mid B] = 20\,\mathrm{ms},\, \operatorname{Var}(\text{delay}\ \mid B) = 25$
\end{itemize}

Find (a) $E[\text{delay}]$, (b) $\operatorname{Var}(\text{delay})$.

\PSsolution{} 

\begin{enumerate}
\def\labelenumi{(\alph{enumi})}
\tightlist
\item
  $E[D] = 0.7(5) + 0.3(20) = 3.5 + 6 = 9.5$ ms
\item
  $E[\operatorname{Var}(D \mid \text{type})] = 0.7(2) + 0.3(25) = 1.4 + 7.5 = 8.9$\PSbr{}$\operatorname{Var}(E[D \mid \text{type}]) = 0.7(5^{2}) + 0.3(20^{2}) - 9.5^{2} = 17.5 + 120 - 90.25 = 47.25$\PSbr{}$\operatorname{Var}(D) = 8.9 + 47.25 = 56.15$
\end{enumerate}

\end{PSbox}

\begin{PSbox}[unbreakable]{prob}{Problem 2}

A random number $N$ of components are tested. $N \sim \mathrm{Poisson}(10)$. Each has independent failure probability $p = 0.05$. Let $X$ = number of failures. Find $E[X]$ and $\operatorname{Var}(X)$.

\PSsolution{} $X \mid N \sim \mathrm{Binomial}(N,\, 0.05)$. $E[X \mid N] = 0.05\,\mathrm{N}$. $\operatorname{Var}(X \mid N) = N(0.05)(0.95) = 0.0475\,\mathrm{N}$.\PSbr{}$E[X] = E[E[X \mid N]] = E[0.05\,\mathrm{N}] = 0.05(10) = 0.5$\PSbr{}$\operatorname{Var}(X) = E[\operatorname{Var}(X \mid N)] + \operatorname{Var}(E[X \mid N]) = E[0.0475\,\mathrm{N}] + \operatorname{Var}(0.05\,\mathrm{N}) = 0.0475(10) + 0.05^{2}(10) = 0.475 + 0.025 = 0.5$

\end{PSbox}

\begin{PSbox}[unbreakable]{prob}{Problem 3}

Write MATLAB code to verify Problem 2 by simulation.

\PSsolution{} 

\tcbset{PScodemode/.style={unbreakable}}

\begin{Shaded}
\begin{Highlighting}[]
\VariableTok{N\_trials} \OperatorTok{=} \FloatTok{100000}\OperatorTok{;} \VariableTok{p} \OperatorTok{=} \FloatTok{0.05}\OperatorTok{;} \VariableTok{lambda} \OperatorTok{=} \FloatTok{10}\OperatorTok{;}
\VariableTok{N} \OperatorTok{=} \VariableTok{poissrnd}\NormalTok{(}\VariableTok{lambda}\OperatorTok{,} \FloatTok{1}\OperatorTok{,} \VariableTok{N\_trials}\NormalTok{)}\OperatorTok{;}
\VariableTok{X} \OperatorTok{=} \VariableTok{binornd}\NormalTok{(}\VariableTok{N}\OperatorTok{,} \VariableTok{p}\NormalTok{)}\OperatorTok{;}
\VariableTok{fprintf}\NormalTok{(}\SpecialStringTok{\textquotesingle{}E[X]: theory=0.5, sim=\%.3f\textbackslash{}n\textquotesingle{}}\OperatorTok{,} \VariableTok{mean}\NormalTok{(}\VariableTok{X}\NormalTok{))}\OperatorTok{;}
\VariableTok{fprintf}\NormalTok{(}\SpecialStringTok{\textquotesingle{}Var(X): theory=0.5, sim=\%.3f\textbackslash{}n\textquotesingle{}}\OperatorTok{,} \VariableTok{var}\NormalTok{(}\VariableTok{X}\NormalTok{))}\OperatorTok{;}
\end{Highlighting}
\end{Shaded}

\emph{Next Module: {Module 9 — Covariance and Correlation}}

\end{PSbox}
